%% Master CV - Enxu Liu
%% Comprehensive overview of all competencies, experience, and achievements.
%% Use as a master reference when tailoring role-specific CVs.
%%
%% NOTE: This master intentionally exceeds the 2-page limit that applies to
%% tailored CVs. It is a source to cut FROM, not a document to submit.
%% Tailored CVs (cv/main_<company>_<role>.tex) must be exactly 2 pages.
%%
%% Compile with: pdflatex main_example.tex   (run twice for page refs)
%%
%% ENGINE NOTE for this machine (Debian TeX Live 2019, moderncv 2.0.0):
%% use pdflatex, NOT lualatex. lualatex routes through fontspec, and this
%% install's luaotfload font-name database cannot resolve OpenType fonts by
%% name (Latin Modern and FontAwesome both fail). pdflatex uses Type 1 fonts
%% and compiles cleanly. The lualatex advice in 05-cv-templates.md applies to
%% modern MiKTeX with fontawesome5, which is a different stack.

\documentclass[11pt,a4paper,sans]{moderncv}
\moderncvstyle{banking}
\moderncvcolor{blue}

% Force the name AND section headings to render in moderncv blue (color1).
% Defaults leave these black, which clashes with the rest of the blue accent scheme.
%
% moderncv >= 2.1 splits the name into \firstnamestyle and \lastnamestyle;
% moderncv 2.0.0 (shipped with TeX Live 2019 / Debian) has only \namestyle.
% Detect which is available so this file compiles on both.
\makeatletter
\@ifundefined{firstnamestyle}{%
  \renewcommand*{\namestyle}[1]{{\fontsize{34}{36}\bfseries\upshape\color{color1}#1}}%
}{%
  \renewcommand*{\firstnamestyle}[1]{{\fontsize{34}{36}\bfseries\upshape\color{color1}#1}}%
  \renewcommand*{\lastnamestyle}[1]{{\fontsize{34}{36}\bfseries\upshape\color{color1}#1}}%
}
\makeatother
\renewcommand*{\sectionstyle}[1]{{\sectionfont\color{color1}#1}}

\usepackage[utf8]{inputenc}
\usepackage{needspace}
% NOTE: do NOT \usepackage{hyperref} here - moderncv.cls already loads it with
% the [unicode] option, and re-loading triggers an option clash.
%
% moderncv loads hyperref inside \AtEndPreamble, so \hypersetup does not exist
% yet at this point in the preamble. Our settings must be deferred into the same
% hook, which runs after the class has loaded hyperref.
\AtEndPreamble{\hypersetup{
    colorlinks=true,
    linkcolor=blue,
    filecolor=magenta,
    urlcolor=blue,
    pdftitle={Enxu Liu - CV},
}}
\usepackage[scale=0.80]{geometry}
\usepackage{import}

% personal data
\name{Enxu}{Liu}
\address{Ann Arbor, Michigan, USA}{}{}
\phone[mobile]{(734) 489-3536}
\email{enxu@umich.edu}
\extrainfo{\href{https://github.com/Liu-Enxu}{github.com/Liu-Enxu}}

\begin{document}

\makecvtitle

% ============================================================
%     PROFILE STATEMENT
% ============================================================

\vspace{6pt}
\small{Embedded systems engineer with an MSc in Electrical \& Computer Engineering from the University of Michigan, working across firmware, embedded Linux, and hardware bring-up. Owned a consumer product end to end at a startup, from FreeRTOS architecture and PCB design through bring-up, FCC compliance, and shipping to sale. Equally at home writing a custom FreeRTOS allocator to fit three middlewares into 128\,KB of SRAM and taking a Rockchip SoC from UART console through U-Boot and kernel to a working Linux shell. Co-inventor on three provisional patents arising from the underlying research.}

% ============================================================
%     CORE COMPETENCIES
% ============================================================

\section{Core Competencies}
\vspace{1pt}
\begin{itemize}

\item \textbf{Embedded firmware}: Embedded C/C++, FreeRTOS (tasks, timers, heap\_4 internals, custom allocators), bare-metal BSP design, bootloaders, delta OTA with rollback, LVGL, FatFS.

\item \textbf{Firmware security}: eFuse hardware root of trust; RSA-signed firmware images with a verified boot chain and secure bootloader; anti-rollback protection against downgrade; BLE application-level handshake and Resolvable Private Address (RPA) privacy; secure debug lockout on production units; MISRA C on safety-relevant code.

\item \textbf{Embedded Linux \& BSP}: Yocto (scarthgap, bbappend authoring, BitBake), Buildroot, mainline U-Boot (SPL, DDR init), kernel build and cross-compilation, custom kernel modules, device tree authoring and pin-mux debugging, V4L2, MIPI-CSI and Rockchip ISP, RKNN NPU inference.

\item \textbf{MCU \& SoC platforms}: ESP32 (Xtensa LX6) and ESP32-C3 (RISC-V) with ESP-IDF, STM32F446RE (ARM Cortex-M4, CubeMX/HAL), NXP S32K144 (ARM Cortex-M4F), Raspberry Pi 4/5, Orange Pi 5 (RK3588S), Rockchip RV1126, Allwinner V3s.

\item \textbf{Buses \& peripherals}: I2C, SPI, UART, CAN / CAN-FD, BLE, TCP/IP, USB gadget, ADC, DAC, PWM, DMA, parallel 8080.

\item \textbf{Controls \& robotics}: ROS2 (node architecture, topics, custom IPC), PID, MPC, LQR/LQE/LTR, Kalman filtering and EKF, GNSS/INS sensor fusion, BLDC-FOC motor control, MATLAB/Simulink/Stateflow.

\item \textbf{Hardware \& lab}: PCB design (KiCAD), LTspice simulation, oscilloscope, logic analyzer, multimeter, JTAG/SWD (J-Link, ST-Link), board bring-up and root-cause analysis, DFM/DFA, FCC compliance.

\item \textbf{Tooling \& process}: Git/GitHub, CMake, scons, GDB, Keil, VS Code, Python, Qt (C++), bash, MISRA C, Jira, Scrum, CI/CD.

\end{itemize}

% ============================================================
%     PROFESSIONAL EXPERIENCE
% ============================================================

\section{Professional Experience}
\vspace{3pt}
\begin{itemize}

\needspace{5\baselineskip}
\item{\cventry{2024--Present}{Mechatronics Engineer}{Motion Sync (spinoff startup)}{Ann Arbor, MI}{}{\vspace{1pt}
\textit{Haptic Pillow -- consumer product, owned firmware and hardware end to end}
\begin{itemize}
    \item Led FreeRTOS firmware design for ESP32 (ESP-IDF) with fully decoupled peripheral drivers; the architecture proved migratable across hardware. Built with CMake.
    \item Built an object-oriented BSP: BMI160 IMU (I2C), W25Q NOR flash (SPI), vibration motor (PWM) at millisecond accuracy via FreeRTOS timers.
    \item Designed a delta OTA pipeline over BLE on an A/B dual-bank boot partition with fail-safe rollback and anti-rollback protection: patch size reduced to 1.2\,KB, update completing in under 10\,s. Optimized the iOS app link for slave latency and throughput.
    \item Implemented battery management: single-cell Li-ion monitoring (ADC), power-path/MOSFET control, battery-sag debugging; low-power configuration cut BLE TX power by ${\sim}$1.94\,mW.
    \item Designed and brought up the PCB (ESP32-C3, IMU, DRV8833, TP4056) from datasheet and schematic with LTspice simulation; root-caused an intermittent brown-out to supply-rail sag under motor inrush using an oscilloscope and logic analyzer.
    \item Built the device security chain: an eFuse hardware root of trust anchoring a verified boot chain, RSA-signed images checked by a secure bootloader before execution, anti-rollback protection against downgrade, a BLE application-level handshake, RPA addressing for privacy, and secure debug lockout on production units.
    \item Carried the product end to end to sale, shipping 200 units; DFM/DFA reduced cost by 30\%. Agile/Scrum, Jira, Git/GitHub CI/CD.
\end{itemize}
\vspace{2pt}
\textit{Active seat programme (PREACT) -- research through Tier 1 automotive OEM delivery}
\begin{itemize}
    % CONDITIONAL - include only when the posting mentions research, R&D, spinoff, clinical or
    % human-subjects work. Cut by default on firmware/embedded CVs. See 01-candidate-profile.md.
    \item Delivered an actuation and data-acquisition system that reduced motion-sickness accumulation by \textbf{${\sim}$19\%} across a 27-subject vehicle study, verified by a linear mixed-model fit (slope difference $p=1.8\mathrm{e}{-5}$) and corroborated by EDA physiological data.
    \item Architected a ROS2 DAQ on Raspberry Pi 5 with a custom Yocto (scarthgap) image, gcc cross-compiled and BitBake-flashed; one sensor per node, core-pinned for real-time jitter isolation; DAQ GUI in Qt (C++).
    \item Deployed the system onto a local Tier 1 automotive OEM seat platform with tunable timing and roll/pitch motion profiles, tapping the OEM vehicle CAN bus for additional sensing and upgrading triggering to computer-vision-based prediction.
    \item Authored a bbappend adding \texttt{KERNEL\_MODULE\_AUTOLOAD:rpi += "i2c-dev"}; validated I2C/UART/BLE/TCP on hardware. Custom IPC via Unix-domain sockets and multiprocessing.
    \item Time-synchronized IMU (I2C), video (ffmpeg/V4L2), dual GNSS (on-board u-blox plus TCP-networked), and BLE into ROS2 topics with ns-precision timestamps and auto-reconnect.
    \item Delivered BLDC motor control with field-oriented control (FOC) on real hardware, commanded over CAN-FD: integrated the CAN-FD setpoint and telemetry path and implemented portions of the FOC control loop.
    \item Deployed a Kalman filter for lane-keep-assist ADAS and a PID controller for seat actuation. Data processing in Python, datasets on AWS.
\end{itemize}
\vspace{2pt}
\textit{Automotive vision prediction -- openpilot port to Orange Pi 5 (RK3588S)}
\begin{itemize}
    \item Ported sunnypilot/openpilot to Orange Pi 5 on Ubuntu-rockchip; built with scons, resolving Git-LFS model-weight and Mali OpenCL toolchain issues.
    \item Ran the supercombo driving model in fp32 on the Mali-G610 GPU (OpenCL/tinygrad) at ${\sim}$10\,Hz for live path, lane, and lead-vehicle prediction on edge hardware.
    \item Brought up USB camera capture (V4L2, MJPG) into a VisionIPC shared-memory pipeline feeding the model and a Qt UI overlay; calibrated camera intrinsics with OpenCV.
    \item Removed the vehicle-CAN dependency by fusing GNSS (NEO-M9N) and IMU (BMI160) in an Innovation-Adaptive-Estimation Kalman Filter with IGGIII robust weighting, measurably more robust than a vanilla KF under measurement noise and GNSS dropout.
\end{itemize}}}

\end{itemize}

% ============================================================
%     SELECTED PROJECTS
% ============================================================

\section{Selected Projects}
\vspace{3pt}
\begin{itemize}

\needspace{5\baselineskip}
\item{\cventry{2026--Present}{MIPI-CSI camera and embedded Linux BSP}{RV1126 Edge Camera}{}{}{\vspace{1pt}
\begin{itemize}
    \item Brought up the ALIENTEK RV1126 (Rockchip) board: UART console, U-Boot to Linux 4.19 kernel to shell; built the full Buildroot vendor SDK image and flashed via rkdeveloptool.
    \item MIPI-CSI IMX335 camera bring-up through the rkcif / rkisp / rkispp ISP pipeline; captured H.264-encoded 2K video with live LCD preview.
    \item Modified the vendor SDK device tree (GPIO/LED node), rebuilt kernel/DTB, verified in \texttt{/sys/class/leds} and \texttt{/proc/device-tree}.
    \item Loaded the NPU (galcore) for on-device inference; V4L2 capture with RKNN YOLO detection overlay in development.
\end{itemize}}}

\vspace{3pt}

\needspace{5\baselineskip}
\item{\cventry{2026}{Personal project}{Embedded Linux Bring-Up: Lichee Pi Zero (V3s)}{}{}{\vspace{1pt}
\begin{itemize}
    \item Cross-compiled mainline U-Boot (SPL to DDR2 init) and the Linux kernel from source; flashed SPL to the sunxi 8\,KB offset and booted to a shell over UART.
    \item Brought the board up headless: serial console, SSH over USB-gadget/Ethernet, expanded ext4 rootfs, deployed zImage and device tree to the boot partition.
    \item Wrote a custom kernel module driving an onboard GPIO LED (gpiolib, kernel timer, module params), verified via insmod/rmmod and dmesg.
    \item Resolved a device-tree pin-mux conflict by disabling SDIO mmc1 to reclaim the LED GPIOs; enabled the DWMAC\_SUN8I Ethernet driver and configured DDR2/LCD via Kconfig.
\end{itemize}}}

\vspace{3pt}

\needspace{5\baselineskip}
\item{\cventry{2026--Present}{Personal project}{STM32 Linux-Inspired Handheld System}{}{}{\vspace{1pt}
\begin{itemize}
    \item Designed HAL firmware on STM32F446RE (ARM Cortex-M4) integrating FreeRTOS, LVGL, and FatFS through a single ucHeap; gcc and Keil toolchains, ST-Link and gdb debug.
    \item Implemented a custom \texttt{pvPortRealloc} over the FreeRTOS heap\_4 allocator, handling the split/coalesce edge cases the upstream API omits; booted the full stack within 128\,KB SRAM and 512\,KB flash.
    \item Built a bare-metal BSP and object-oriented LVGL GUI on an ILI9486 LCD with ADC+DMA touch sampling over a parallel 8080 interface; debugged a stale-coordinate fault to a non-clearing ADC conversion-complete flag using a logic analyzer.
\end{itemize}}}

\vspace{3pt}

\needspace{5\baselineskip}
\item{\cventry{2024}{Graduate coursework}{Motor Embedded Motion Control}{}{}{\vspace{1pt}
\begin{itemize}
    \item Configured PWM, ADCs, and interrupt routines in C on an NXP S32K144 to control BLDC motor torque.
    \item Defined CAN-FD message IDs for steering-angle and torque setpoints with a Stateflow state machine for mode arbitration; tuned a PID loop to under 10\% overshoot, HIL-verified in real time.
\end{itemize}}}

\vspace{3pt}

\needspace{4\baselineskip}
\item{\cventry{2024}{Graduate coursework}{MIMO Feedback Controller Design}{}{}{\vspace{1pt}
\begin{itemize}
    \item Developed PID controllers via loop shaping for a simulated wafer-etching system in MATLAB/Simulink.
    \item Designed a MIMO feedback controller integrating LQR with integrators, LQE, and loop transfer recovery; verified robustness by margin analysis.
\end{itemize}}}

\vspace{3pt}

\needspace{4\baselineskip}
\item{\cventry{2024}{Graduate coursework}{Self-Driving Cars: Path Planning and Motion Control}{}{}{\vspace{1pt}
\begin{itemize}
    \item Developed reinforcement-learning motion planning (Markov Decision Process) with MPC motion control, simulated and verified in ROS Rviz.
\end{itemize}}}

\end{itemize}

% ============================================================
%     EDUCATION
% ============================================================

\section{Education}
\vspace{1pt}
\begin{itemize}

\needspace{4\baselineskip}
\item{\cventry{2023--2024}{MSc in Electrical \& Computer Engineering}{University of Michigan}{Ann Arbor, MI}{\textit{GPA 3.97}}{\vspace{1pt}
Embedded systems, controls, and robotics.
}}

\vspace{3pt}

\needspace{4\baselineskip}
\item{\cventry{2020--2023}{BEng in Electronic \& Electrical Engineering}{University College London}{London, UK}{\textit{First-Class Honours, GPA 4.00}}{\vspace{1pt}
Electronics, signals, and control systems.
}}

\end{itemize}

% ============================================================
%     PATENTS
% ============================================================

\section{Patents}
\vspace{1pt}
\begin{itemize}
\item{Co-inventor, three University of Michigan provisional patent applications (2026), UM Ref. Nos. 2026-235-01, 2026-234-01, 2026-233-01:
\begin{itemize}
    \item \textit{Systems and Methods Implementing Control Logic Software Architecture for Preemptive Control of Active Devices and Systems Within a Vehicle}
    \item \textit{Systems and Methods for Preemptive Control of Active Notification Devices and Systems Within a Vehicle}
    \item \textit{Systems and Methods Implementing Perception and Prediction Multithreaded Architecture for Preemptive Control of Active Devices and Systems Within a Vehicle}
\end{itemize}}
\end{itemize}


\end{document}
